\section{Obstructions, finite rings and partial degrees}\label{sec:line-obstructions}
\begin{theorem}[Necessary data on the line]\label{thm:line-obstructions}
In both line categories, $E$ forces $L$ false, without any uniform
attraction assumption. Multiple TI stationary states force $E$ false
and $L$ true. Always $N_{\Br}\le N_{\ord}\le N$; if $N=1$, then
$N_{\ord}=1$ exactly when $L$ holds.
The total count belongs to $\{1,2,\ldots,\aleph_0,\frakc\}$.
Consequently the only possible coarse data $(N,E,L)$ are
\[
 (1,\mathrm{true},\mathrm{false}),\quad
 (1,\mathrm{false},\mathrm{false}),\quad
 (1,\mathrm{false},\mathrm{true}),\quad
 (k,\mathrm{false},\mathrm{true}),
\]
where the last case has $k\in\{2,3,\ldots,\aleph_0,\frakc\}$.
At \eqref{eq:line-rate}, the first case and the last with $k=\frakc$
are realized in both categories. No realization or exclusion of
the other cases is established here.
\end{theorem}
\begin{proof}
Nonemptiness and compactness follow from finite marginal coordinates,
closed TI and stationarity equations, and Ces\`aro averaging.
For distinct TI stationary states choose a real local, or Hermitian,
observable with means $a\ne b$. In their half-mixture the variance
of its normalized length-$l$ spatial average is at least
$(a-b)^2/4$ for every $l$. If all distant covariances decayed, a
near/far split of its double sum would make that variance tend to
zero: only $O(l)$ pairs have any fixed bounded displacement.
Thus multiplicity gives $L$ and precludes $E$. The Bragg/covariance
implication proved in Lemma~\ref{lem:two-order} gives the count
inequalities. If $N=1$, the compact convex TI stationary set is a
singleton, so its only point is ordered exactly when $L$ holds.

We prove that $E$ forbids nondecaying covariance. In the classical
case, let $\mu$ be the unique attracting law and $P$ the transition
operator on observables. For each deterministic initial row $\eta$,
$P^tf(\eta)\to\mu(f)$. Bounded convergence gives convergence to this
constant in $L^2(\mu)$. For distant enough translations the two
$t$-step backward cones are disjoint. Independent innovations and
stationarity imply
\[
 \Cov_\mu(f,\tau_zg)
   =\Cov_\mu(P^tf,\tau_zP^tg).
\]
Cauchy--Schwarz bounds the right side by the product of the centered
$L^2$ norms, uniformly in sufficiently large $z$. Letting $t\to\infty$
proves covariance decay without uniform convergence over starts.

For the quantum case write $P$ for the Heisenberg round and $\mu$
for its unique stationary state, which is TI by uniqueness.
Suppose a local covariance fails to decay and choose $|z_j|\to\infty$
along which it converges to a value different from the product of
means. In the cyclic GNS representation $(\pi,\Omega)$ of $\mu$,
the Hilbert space is separable: the chain algebra has a countable
norm-dense subset. By a diagonal subsequence the bounded operators
$\pi(\tau_{z_j}b)$ converge weakly simultaneously for a countable
dense set of $b$, hence for all $b$ by norm approximation.
Call the limit $Z(b)$. It belongs to the generated von Neumann
algebra and commutes with every local operator, by eventual spatial
disjointness and local approximation. Thus $Z$ is positive, unital
and linear into the center. Define
\[
 \Lambda(a,b)=\langle\Omega,\pi(a)Z(b)\Omega\rangle.
\]
For sufficiently distant local $a,b$, independent backward circuits
give $P(a\tau_zb)=P(a)\tau_zP(b)$. Stationarity and translation
covariance, then the weak limit, imply
$\Lambda(a,b)=\Lambda(Pa,Pb)$, and by iteration the same identity
for every $t$. The chosen covariance means that $\Lambda(f,g)$
is not $\mu(f)\mu(g)$ for some local pair.

We now express $\Lambda$ as an integral of products of two states.
For finitely many Hermitian $b$, partition the spectra of their
commuting central $Z(b)$ into intervals of width $\eta$, and let
$Q_i$ be the joint central spectral projections. For
$p_i=\langle\Omega,Q_i\Omega\rangle>0$ put
\[
 \alpha_i(a)=\frac{\langle\Omega,Q_i\pi(a)\Omega\rangle}{p_i},
 \qquad
 \beta_i(b)=\frac{\langle\Omega,Q_iZ(b)\Omega\rangle}{p_i}.
\]
Centrality and positivity make these positive unital functionals
on the chain algebra. On each projection, $Z(b)$ differs from its
conditional scalar $\beta_i(b)$ by norm at most $\eta$ for the
chosen $b$. Consequently
\[
 \left|\Lambda(a,b)-\sum_i p_i\alpha_i(a)\beta_i(b)\right|
 \le\eta\|a\|.
\]
Exhaust a countable dense Hermitian set and let $\eta\downarrow0$.
The product of the two compact metrizable chain state spaces has
compact probability-measure space, so a subsequence of the laws
$\sum_i p_i\delta_{(\alpha_i,\beta_i)}$ converges to a probability
law $\theta$. It gives
$\Lambda(a,b)=\int\alpha(a)\beta(b)\,d\theta$ for every $a,b$.
For each pair of states, pointwise attraction gives
$\alpha(P^ta)\to\mu(a)$ and $\beta(P^tb)\to\mu(b)$.
Their product is bounded by $\|a\|\|b\|$. Dominated convergence,
together with invariance of $\Lambda$, forces
$\Lambda(a,b)=\mu(a)\mu(b)$, a contradiction.
Thus $E$ implies $L$ false in both categories.

The nonextreme points of a compact metrizable convex set form a
countable union of compact sets of midpoints of pairs at distance
at least $1/j$. Hence its extreme set is $G_\delta$ and is Polish.
The perfect-set theorem gives the asserted possibilities for $N$,
without a continuum-hypothesis assumption. The implications already
proved give the displayed necessary coarse list.
A constant classical table gives its one-step iid replacement law.
A fixed on-site quantum preparation followed by qubit noise gives
a product stationary state. Both have
$(1,0,0,\mathrm{true},\mathrm{false})$; a one-symbol classical
alphabet is also allowed. A fixed never-halting controller gives
$\sigma_\infty$ by Proposition~\ref{prop:line-phases}.
These are the two claimed realizations. The Rice reductions range
over the actual image and require no classification of the remaining
possibilities.
\end{proof}

\begin{remark}[Fixed finite rings]\label{rem:line-rings}
A classical ring of $l$ sites has a complete reset with probability
$e^l$ each step. Its whole-state total-variation contraction
coefficient is at most $1-e^l$. If a finite quantum ring has $k$
physical wires, simultaneous reset of every wire in the final layer
has probability $e^k$ and makes the full output trace-and-prepare.
Its trace-distance contraction coefficient is at most $1-e^k$.
Iteration gives a unique stationary state attracting every initial
state in either case. An unbounded-distance covariance label concerns
the infinite line and has no such meaning on one fixed ring. These
contraction bounds need not be uniform as the ring grows; growing
storage horizons are consistent with fixed-ring uniqueness.
At replacement probability one every input is reset; the theorems
use the strictly smaller rate \eqref{eq:line-rate}. A noisy identity
also has the reset datum, since a fixed finite observation forgets
its starting values after all its sites or wires have been reset.
\end{remark}

\begin{theorem}[Partial degree results on the line]\label{thm:line-cuts}
Uniformly on valid presentations in either line category, for each
$r\ge1$ the cut $N\le r$ belongs to $\Pi^0_2$ and is halting-hard.
For each $k\ge2$ the cut $N\ge k$ belongs to $\Sigma^0_2$ and is
nonhalting-hard. Matching second-level hardness on the line is not
asserted.
\end{theorem}
\begin{proof}
We give the effective finite-coordinate argument, including its
hypotheses for the per-layer channel. Enumerate real local coordinates
$x_1,x_2,\ldots$ bounded by one: cylinder probabilities classically,
and real and imaginary density-matrix entries with Hermiticity
built in quantumly. Use all finite intervals and their translates.
Normalization, consistency, TI and stationarity are countably many
effective linear constraints. Classical positivity is nonnegativity.
Quantum positivity is the family
$v^*\rho_Iv\ge0$ for every rational complex vector $v$ on every
finite interval $I$; density makes this equivalent to positive
semidefiniteness. The backward circuit computing a local Heisenberg
image includes every depolarizing layer. Rational $e$ and algebraic
gates give exact algebraic coefficients with decidable arithmetic
and order. Thus all these linear constraints are effectively
enumerable in both categories.

Let $b(m)\ge m$ increase to cover the coordinates used by the first
$m$ constraints. Let $P_m$ be their nonempty bounded algebraic
polytope in those coordinates with coordinate bounds $[-1,1]$.
All earlier constraints are retained. Every genuine stationary
state projects into every $P_m$. Conversely a coordinatewise limit
of outer points with $m\to\infty$ satisfies every constraint and
is genuine. Product compactness implies that for each fixed $n$
the first-$n$ projections of $P_m$ converge in Hausdorff distance
to those of $\cK$: otherwise distant outer points have a convergent
subsequence contradicting that characterization.

Every vertex of $P_m$ is computable: enumerate full-rank active
constraint systems, solve, and check all inequalities. This also
handles lower-dimensional polytopes, since at a vertex active
normals must span the ambient space or a feasible segment remains.
A bounded polytope is the convex hull of its vertices, by induction
on face dimension using boundary-to-boundary segments. Exact
algebraic linear feasibility is therefore decidable.
Put $\eta=2^{-n}$. If $V$ is the number of vertices, choose an integer
$L_{m,n}\ge8V/\eta$ and let $\mathcal Q_{m,n}$ be all their convex
combinations whose coefficients are multiples of $1/L_{m,n}$.
This finite set is an $\eta/4$-net in the first $n$ coordinates:
round the first $V-1$ weights down and put the remainder on the
last vertex; total weight error is at most $2(V-1)/L_{m,n}$.

For $m\ge n$, let $\mathsf T_r(n,m)$ mean that some $r$ points of
$\mathcal Q_{m,n}$, with repetitions, have a convex hull within
$\eta$ of every vertex in first-$n$ maximum norm. For each candidate
and vertex this is finite algebraic feasibility in simplex
coefficients $\lambda_i$, with inequalities
$|v_j-\sum_i\lambda_i y_{i,j}|\le\eta$, $1\le j\le n$.
So $\mathsf T_r$ is decidable. Vertex coverage covers the polytope
because the sum of a convex hull and a norm ball is convex.
We claim
\[
 N\le r\quad\Longleftrightarrow\quad
 \forall n\ge1\ \exists m\ge n:\mathsf T_r(n,m).
\]
A compact convex set with finitely many extremes is their hull:
if a point lay outside, a finite-coordinate separating functional
would have a maximizing compact face disjoint from that hull;
successive coordinate maximization there gives a further extreme,
a contradiction. Hence if $N\le r$, choose $r$ generators, repeating
as needed. For fixed $n$, choose $m$ so the outer projection is within
$\eta/4$ of $\cK$, then approximate those generators by
$\mathcal Q_{m,n}$ within $\eta/4$. Their hull covers the outer
projection within $\eta/2$, proving the test.

Conversely choose successful stages $m_n\ge n$ and their $r$ centers.
A common coordinatewise convergent subsequence has limits
$y_1,\ldots,y_r\in\cK$, since every constraint eventually holds.
For any $x\in\cK$, vertex coverage provides simplex coefficients
approximating its first $n$ coordinates within $2^{-n}$. A further
subsequence of these coefficients converges, giving
$x=\sum_i\lambda_i y_i$ in each coordinate. The limiting centers
were chosen before $x$, so these same $r$ points generate all of
$\cK$. There are at most $r$ extremes. This proves the equivalence
and the $\Pi^0_2$ membership; complementation gives the $\Sigma^0_2$
membership of $N\ge k$.

Use a reset host in Proposition~\ref{prop:line-branches}. Its two
counts are $1$ on halting and $\frakc$ on nonhalting. This reduces
halting to every listed finite upper cut and nonhalting to every
listed finite lower cut. Such a two-branch reduction gives no
matching hardness at the second arithmetical level. The totality
construction in Section~\ref{sec:completeness} uses transverse
slices and a different controller and is not a line construction.
\end{proof}
